\chapter{Grassmanniany i współrzędne Plückera}
\label{ch:lin29-grassmanniany}
\index{grassmannian}
\index{współrzędne Plückera}
\index{wiązka tautologiczna}

W Rozdziałach~\ref{ch:lin10-iloczyny-tensorowe}
i~\ref{ch:lin11-algebry} zbudowaliśmy iloczyn tensorowy
i algebrę zewnętrzną. Obecnie zbadamy zbiór wszystkich
podprzestrzeni ustalonego wymiaru.
Skonstruujemy atlas grassmannianu, jego przestrzeń styczną
i wiązkę tautologiczną, a następnie przedstawimy
współrzędne Plückera oraz odwzorowanie Gaussa.
Przykłady w małych wymiarach zilustrują konstrukcje.

\section{Grassmannian jako przestrzeń podprzestrzeni}
\label{sec:lin29-grass-atlas}

W całym rozdziale przez $\mathbb K$ oznaczamy
$\R$ albo $\C$.
Jeśli $0\leq k\leq n$, \emph{grassmannian}
$\operatorname{Gr}_k(\mathbb K^n)$ to zbiór wszystkich
$k$-wymiarowych podprzestrzeni liniowych $\mathbb K^n$.
Jego \emph{punktem} jest cała podprzestrzeń $L$,
a nie wektor w $L$. Na przykład
$\operatorname{Gr}_1(\R^n)=\mathbb{RP}^{n-1}$
, gdzie $\mathbb{RP}^{n-1}$ oznacza przestrzeń prostych rzeczywistych, a
$\operatorname{Gr}_1(\C^n)=\mathbb{CP}^{n-1}$.
Dla $k=0$ i $k=n$ grassmannian jest jednym punktem.
Dla dowolnej $n$-wymiarowej przestrzeni $V$ nad $\mathbb K$
piszemy analogicznie $\operatorname{Gr}_k(V)$; wybór bazy
identyfikuje go z powyższym modelem, lecz konstrukcje nie zależą od tej bazy.

\begin{twierdzenie}[Mapy wykresowe grassmannianu]
\label{thm:lin29-grass-atlas}
Niech $\mathbb K^n=P\oplus Q$, gdzie $\dim_{\mathbb K}P=k$.
Zbiór
\[
 U_P=\{L\in\operatorname{Gr}_k(\mathbb K^n):
           \pi_P|_L:L\longrightarrow P\text{ jest izomorfizmem}\}
\]
można jednoznacznie opisać odwzorowaniami liniowymi
$A:P\to Q$: punktowi $A$ odpowiada jego
\emph{wykres} $\Gamma(A)=\{u+Au:u\in P\}$.
Mapy $\Gamma(A)\mapsto A$ tworzą gładki atlas
grassmannianu. Jego wymiar nad $\mathbb K$ wynosi
$k(n-k)$. W przypadku zespolonym atlas jest holomorficzny,
a rzeczywisty wymiar wynosi $2k(n-k)$.
Oznaczenie $U_P$ zależy również od wybranego dopełnienia $Q$;
pomijamy je w indeksie tylko dla skrócenia zapisu.
\end{twierdzenie}
\begin{proof}
Jeśli $L=\Gamma(A)$, to $\pi_P(u+Au)=u$, więc
$\pi_P|_L$ jest izomorfizmem. W drugą stronę, dla
$L\in U_P$ określmy
\[
 A_L:=\pi_Q|_L\circ(\pi_P|_L)^{-1}:P\longrightarrow Q.
\]
Każdy $v\in L$ ma postać $u+\pi_Qv=u+A_Lu$
dla $u=\pi_Pv$. Stąd $L=\Gamma(A_L)$; dowiedliśmy
zarówno istnienia, jak i jednoznaczności parametru.

Sprawdźmy przejścia między mapami, a nie tylko ich
liczbę. Niech $\mathbb K^n=P'\oplus Q'$ będzie drugim
rozkładem. Dla $A:P\to Q$ połóżmy
$T_A:P\to\mathbb K^n$, $T_A(u)=u+Au$.
Wykres $\Gamma(A)$ należy do $U_{P'}$ dokładnie wtedy,
gdy $D_A:=\pi_{P'}T_A:P\to P'$ jest odwracalny.
Po wybraniu baz $P,P'$ wyznacznik macierzy $D_A$ zależy
wielomianowo od współczynników $A$;
warunek $\det D_A\ne0$ jest więc otwarty. W nowej mapie
parametr ma postać
\begin{equation}\label{eq:lin29-grass-przejscie}
 A'=\pi_{Q'}T_A\circ(\pi_{P'}T_A)^{-1}:P'\longrightarrow Q'.
\end{equation}
Odwracanie macierzy jest gładkie na zbiorze macierzy
odwracalnych, a nad $\C$ jest tam także holomorficzne
(wzór przez dopełnienia algebraiczne i wyznacznik).
Każde $L$ leży w $U_L$, bo dla rozkładu
$\mathbb K^n=L\oplus Q$ otrzymujemy $A_L=0$.
Atlasy pokrywają zatem cały zbiór. Współczynniki
$A\in\operatorname{Hom}_{\mathbb K}(P,Q)$ tworzą
przestrzeń $\mathbb K^{(n-k)k}$, co daje wymiar.

Definiujemy topologię przez atlas: zbiór jest otwarty, gdy
jego obraz w każdej z map wykresowych jest otwarty.
Otwartość dziedzin przejść i ich ciągłość zapewniają, że
same $U_P$ są otwarte, a mapy są homeomorfizmami.
Uzasadnijmy również pozostałe własności topologiczne.
Wybierzmy standardowy iloczyn skalarny, a nad $\C$ hermitowski;
gwiazdka oznacza operator sprzężony. W tym kroku
ograniczmy się do rozkładów ortogonalnych $Q=P^\perp$;
takie mapy wciąż pokrywają cały grassmannian
(dla danego $L$ bierzemy $P=L$).
Wzór~\eqref{eq:lin29-grass-przejscie} zapewnia zgodność
z pozostałymi mapami. Każdemu $L$
przypiszmy macierz rzutowania ortogonalnego $\Pi_L$.
Jest ona wyznaczona jednoznacznie przez $L$, a jej
współczynniki w mapie wykresowej są gładkie:
\[
 \Pi_{\Gamma(A)}
  =T_A(T_A^*T_A)^{-1}T_A^*.
\]
Wyprowadźmy ten wzór z warunku prostopadłości.
Obraz $T_A$ to $\Gamma(A)$, więc rzut wektora $v$ ma postać
$T_Au$. Różnica $v-T_Au$ ma być prostopadła do wszystkich
$T_Aw$; równoważnie $T_A^*(v-T_Au)=0$.
Stąd $(T_A^*T_A)u=T_A^*v$ i otrzymujemy wypisany wzór.
Macierz $G_A:=T_A^*T_A=I+A^*A$ jest odwracalna, ponieważ
$\langle G_Au,u\rangle=\|u\|^2+\|Au\|^2>0$ dla $u\ne0$.
Odwrót na $U_P$ jest również ciągły w macierzach
rzutowania, ponieważ dla $L\in U_P$
\[
 A_L=\pi_Q\Pi_L|_P\,(\pi_P\Pi_L|_P)^{-1}.
\]
Istotnie, w rozkładzie $P\oplus Q$ mamy
$T_Au=(u,Au)$ i $T_A^*(u,0)=u$.
Zatem $\pi_P\Pi_L|_P=G_A^{-1}$ oraz
$\pi_Q\Pi_L|_P=AG_A^{-1}$, a złożenie w ostatnim wzorze
wynosi $(AG_A^{-1})(G_A^{-1})^{-1}=A$.
Ostatnia odwrotność istnieje, gdyż przy ortogonalnym
rzutowaniu $\pi_P$ odwzorowanie $\Pi_L|_P:P\to L$
jest sprzężeniem izomorfizmu
$\pi_P|_L:L\to P$. Nasza topologia zgadza się więc
z topologią podzbioru macierzy rzutowania: $U_P$ jest w tym
modelu zbiorem, gdzie $\pi_P\Pi_L|_P$ jest odwracalne,
a więc zbiorem otwartym. Oba podane wzory są ciągłe na tych mapach.
Zbiór $\Pi^*=\Pi$, $\Pi^2=\Pi$, $\operatorname{tr}\Pi=k$
jest domknięty i ograniczony w skończenie wymiarowej
przestrzeni macierzy, a więc zwarty. Istotnie, samosprzężone
idempotenty mają wartości własne $0,1$, więc norma operatorowa
jest co najwyżej $1$, a ślad jest rangą. Są zatem dokładnie
rzutami na $k$-płaszczyzny. Ten model jest także
Hausdorffa i ma przeliczalną bazę. Atlas definiuje
zatem rozmaitość w zwykłym sensie gładkiej rozmaitości.
\end{proof}

\begin{figure}[!htbp]
\centering
\begin{tikzpicture}[>=Stealth,line cap=round,line join=round,scale=1.02]
 \filldraw[fill=manifoldblue!11,draw=manifoldblue!70!black,thick]
 (-2.55,-.3)--(1.45,-.3)--(2.55,.3)--(-1.45,.3)--cycle;
 \filldraw[fill=manifoldgreen!14,draw=manifoldgreen!75!black,thick]
 (-2.55,-1.15)--(1.45,.25)--(2.55,1.15)--(-1.45,-.25)--cycle;
 \draw[->,manifoldgold!85!black,very thick] (0,0)--(0,1.75)
  node[above,font=\small] {$Q$};
 \fill[black!75] (0,0) circle (2pt) node[below left,font=\small] {$0$};
 \fill[manifoldblue] (1.675,.15) circle (2pt)
  node[below right,fill=white,inner sep=1pt,font=\small] {$u\in P$};
 \fill[manifoldgreen!80!black] (1.675,.715) circle (2pt)
  node[above right,fill=white,inner sep=1pt,font=\small] {$u+Au\in\Gamma(A)$};
 \draw[->,manifoldred,thick] (1.675,.15)--(1.675,.715)
  node[midway,left,font=\small] {$Au$};
 \node[manifoldblue!70!black,font=\small\bfseries] at (-1.2,.52) {$P$};
 \node[manifoldgreen!65!black,font=\small\bfseries] at (-2.5,-1.48) {$\Gamma(A)$};
\end{tikzpicture}
\caption{Mapa wykresowa grassmannianu. Przestrzeń $P$ jest
punktem odniesienia; pobliska płaszczyzna $\Gamma(A)$ powstaje,
gdy nad $u\in P$ odkładamy przemieszczenie $Au\in Q$.}
\label{fig:lin29-grass-wykres}
\end{figure}

\begin{przyklad}[Mała mapa, a nie cały grassmannian]
\label{prz:lin29-grass-rp1-mapy}
Na $\operatorname{Gr}_1(\R^2)=\mathbb{RP}^1$ ustalmy
$P=\R e_1$ i $Q=\R e_2$. Prosta
$L_a=\operatorname{span}(e_1+ae_2)$ ma współrzędną
$a\in\R$. Brakuje dokładnie jednej prostej, $\R e_2$:
rzutowanie na $P$ jest na niej zerowe. Zatem mapa
wykresowa jest lokalna, podobnie jak mapa stereograficzna
nie pokrywa całej sfery.

W drugiej mapie, z $P'=\R e_2$ i $Q'=\R e_1$,
zapisujemy $L=\operatorname{span}(be_1+e_2)$.
Na części wspólnej obu map $a\ne0$ i
\[
 e_1+ae_2=a\bigl(a^{-1}e_1+e_2\bigr),
 \qquad b=\frac1a.
\]
To najprostszy przypadek wzoru~\eqref{eq:lin29-grass-przejscie}:
druga mapa dopisuje prostą $\R e_2$, której nie ma
w pierwszej; jej współrzędna $b=0$ odpowiada
,,$a=\infty$'' w potocznym opisie.
\end{przyklad}

\section{Przestrzeń styczna i wiązka tautologiczna}
\label{sec:lin29-grass-styczna}

Wektor styczny do grassmannianu opisuje \emph{pierwszy
rząd zmiany podprzestrzeni}. Wektory należące już do
danej podprzestrzeni nie zmieniają jej położenia:
istotna jest część wychodząca poza nią.
Z każdej z $k$ niezależnych osi podprzestrzeni możemy
wyjść w $n-k$ niezależnych kierunkach poprzecznych;
stąd intuicyjnie bierze się wymiar $k(n-k)$.

\begin{twierdzenie}[Styczne jako odwzorowania do ilorazu]
\label{thm:lin29-grass-styczna}
Dla $L\in\operatorname{Gr}_k(V)$ istnieje kanoniczny
izomorfizm
\begin{equation}\label{eq:lin29-grass-styczna}
 T_L\operatorname{Gr}_k(V)
  \cong\operatorname{Hom}_{\mathbb K}(L,V/L).
\end{equation}
Nad $\C$ prawa strona jest zespoloną przestrzenią
styczną; jej podłożem rzeczywistym jest rzeczywista
przestrzeń styczna.
\end{twierdzenie}
\begin{proof}
Niech $L(t)$ będzie gładką krzywą podprzestrzeni,
$L(0)=L$. Dla $v\in L$ wybierzmy gładką krzywą
$v(t)\in L(t)$ z $v(0)=v$. Definiujemy
$D_{L'(0)}(v):=[v'(0)]\in V/L$, gdzie nawias oznacza
klasę modulo $L$. Pokażmy niezależność od wyboru
$v(t)$. Różnica dwóch takich krzywych $w(t)\in L(t)$
ma $w(0)=0$. W gładkiej lokalnej bazie $s_i(t)$
zapisuje się $w(t)=\sum_i a_i(t)s_i(t)$ z
$a_i(0)=0$, zatem
$w'(0)=\sum_i a_i'(0)s_i(0)\in L$. Klasa
$[v'(0)]$ nie zmienia się. Zmiana parametryzacji
krzywej nie jest jedyną możliwą zmianą przedstawiciela wektora
stycznego; niezależność od dowolnej krzywej o tym samym wektorze
stycznym sprawdzimy teraz w mapie.

Ustalmy teraz uzupełnienie $V=L\oplus Q$.
W mapie z Twierdzenia~\ref{thm:lin29-grass-atlas} mamy
$L(t)=\Gamma(A(t))$ i $A(0)=0$.
Wybieramy $v(t)=v+A(t)v$. Wtedy
\[
 D_{L'(0)}(v)=[A'(0)v]\in V/L.
\]
Tak samo $s_i(t)=e_i+A(t)e_i$ dla bazy $(e_i)$ przestrzeni $L$
daje gładką bazę używaną powyżej. Wzór zależy tylko od $A'(0)$,
czyli współrzędnych wektora stycznego, a nie od dalszego przebiegu krzywej.
Jeśli $D_{L'(0)}=0$, to $A'(0)v\in Q\cap L=0$ dla każdego $v$,
więc $A'(0)=0$ i wektor styczny jest zerowy. To dowodzi iniektywności.
Utożsamienie $Q\cong V/L$ pokazuje, że
$D_{L'(0)}$ jest liniowe w $v$. Każdy
$D\in\operatorname{Hom}(L,V/L)$ można otrzymać:
wybierzmy odpowiadające mu odwzorowanie $A:L\to Q$
i krzywą $L(t)=\Gamma(tA)$. Jej pochodna daje
właśnie $D$. Konstrukcja rozpoczęła się od
krzywej i klas modulo $L$, więc otrzymany
izomorfizm nie zależy od wybranego $Q$.
\end{proof}

\begin{przyklad}[Wektor styczny do przestrzeni prostych]
Niech $L_0=\R e_1\in\operatorname{Gr}_1(\R^2)$ i
$L(t)=\R(e_1+t e_2)$. W mapie z
Przykładu~\ref{prz:lin29-grass-rp1-mapy}
współrzędna krzywej to $a(t)=t$, więc jej prędkość
wynosi $a'(0)=1$. Dla wektora $v=s e_1\in L_0$
wybieramy $v(t)=s(e_1+t e_2)\in L(t)$; z definicji
\[
 D_{L'(0)}(s e_1)=[s e_2]\in\R^2/L_0.
\]
Zatem pojedyncza liczba $a'(0)$ jest macierzą
odwzorowania $L_0\to\R^2/L_0$ w bazach
$e_1$ i $[e_2]$. Krzywa
$\widetilde L(t)=\R(\cos t\,e_1+\sin t\,e_2)$
daje ten sam wektor styczny, bo jej współrzędna
$\tan t$ także ma pochodną $1$ w zerze.
\end{przyklad}

\begin{definicja}[Wiązka tautologiczna]
\index{wiazka tautologiczna@Wiązka tautologiczna}
\label{def:lin29-grass-taut}
Nad $\operatorname{Gr}_k(V)$ określamy
\emph{wiązkę tautologiczną}
\[
 \gamma^k
 :=\{(L,v)\in\operatorname{Gr}_k(V)\times V:v\in L\}
 \longrightarrow\operatorname{Gr}_k(V),\qquad (L,v)\mapsto L.
\]
Jej włókno w punkcie $L$ jest dosłownie podprzestrzenią
$L\subseteq V$.
\end{definicja}

\begin{stwierdzenie}[Lokalność i wiązka ilorazowa]
\label{prop:lin29-grass-wiazki}
Wiązka $\gamma^k$ ma rząd $k$. Jest podwiązką
wiązki trywialnej $\operatorname{Gr}_k(V)\times V$.
Istnieje też wiązka ilorazowa $\mathcal Q$ z
włóknem $\mathcal Q_L=V/L$, rzędu $n-k$,
oraz kanoniczny dokładny ciąg wiązek
\begin{equation}\label{eq:lin29-grass-ciag}
 0\longrightarrow\gamma^k
 \longrightarrow\operatorname{Gr}_k(V)\times V
 \longrightarrow\mathcal Q\longrightarrow0.
\end{equation}
Po wybraniu iloczynu skalarnego otrzymujemy
$\mathcal Q\cong(\gamma^k)^\perp$.
\end{stwierdzenie}
\begin{proof}
Na $U_P$ każdy element $L=\Gamma(A)$ zawiera
dokładnie jeden wektor $u+Au$ dla danego $u\in P$.
Odwzorowanie $(L,u)\mapsto(L,u+Au)$ jest zatem lokalną
trywializacją $U_P\times P\cong\gamma^k|_{U_P}$;
jej odwrotność to $(L,v)\mapsto(L,\pi_Pv)$.
Oba odwzorowania są gładkie, a na włóknach liniowe.
W tym samym rozkładzie $V=P\oplus Q$ klasa
$[w]\in V/\Gamma(A)$ ma jednoznacznego
reprezentanta w $Q$: od $w$ odejmujemy
$u+Au\in\Gamma(A)$, gdzie $u=\pi_Pw$.
Stąd lokalnie $\mathcal Q|_{U_P}\cong U_P\times Q$.
Na każdym włóknie inkluzja $L\hookrightarrow V$
jest iniekcją, a naturalne odwzorowanie $V\to V/L$ jest
surjekcją z jądrem $L$, co wyjaśnia \emph{dokładność}
ciągu~\eqref{eq:lin29-grass-ciag}. Iloczyn skalarny
identyfikuje klasę $[v]$ z $(\id-\Pi_L)v\in L^\perp$.
Rzut zależy gładko od $L$, więc jest to izomorfizm wiązek,
a nie tylko rodzina izomorfizmów włókien.
\end{proof}

\begin{uwaga}[Styczna i iloczyn tensorowy]
Twierdzenie~\ref{thm:lin29-grass-styczna} i
izomorfizm tensorowy~\ref{thm:lin10-hom-rzad} dają
\[
 T\operatorname{Gr}_k(V)\cong
       (\gamma^k)^*\otimes\mathcal Q.
\]
Włókno po prawej stronie w $L$ to
$L^*\otimes(V/L)\cong\operatorname{Hom}(L,V/L)$.
Kolejność czynników jest tu ważna: źródłem odwzorowania
jest $L$, a jego wartości należą do $V/L$.
Zgodność tych izomorfizmów z gładkością również widać w mapie:
przy $L=\Gamma(A)$ przyrost $\dot A$ wysyła $u+Au$ na
klasę $\dot A u$. W trywializacji ilorazu
$[w]\mapsto\pi_Qw-A\pi_Pw$ ma on po prostu macierz $\dot A$.
Zatem izomorfizmy włókien sklejają się w gładki izomorfizm wiązek.
\end{uwaga}

\begin{przyklad}[Proste w płaszczyźnie i wstęga Möbiusa]
Na $\operatorname{Gr}_1(\R^2)\cong\mathbb{RP}^1\cong S^1$
wiązka $\gamma^1$ jest wstęgą Möbiusa (co wynika z poniższego obliczenia). Istotnie,
$L(t)=\R(\cos t,\sin t)$, $0\leq t\leq\pi$,
opisuje pełne okrążenie bazy: $L(0)=L(\pi)$.
Wybrany lokalnie wektor jednostkowy
$(\cos t,\sin t)\in L(t)$ wraca jako jego
przeciwieństwo. To właśnie sklejanie włókna
ze zmianą znaku.
W dwóch mapach z Przykładu~\ref{prz:lin29-grass-rp1-mapy} tę zmianę
widać również w funkcji przejścia: jeśli
$\lambda(e_1+ae_2)=\mu(be_1+e_2)$ oraz $b=1/a$,
to porównanie współczynników daje $\mu=a\lambda$.
Na dwóch składowych części wspólnej ($a>0$ i $a<0$)
funkcja przejścia ma przeciwne znaki. Właśnie tego
nie da się usunąć przez jeden globalny wybór
niezerowego wektora w każdym włóknie: zmiana
lokalnej ramki (czyli lokalnego nigdzie niezerowego przekroju)
przez funkcję nigdzie niezerową
ma na każdej spójnej mapie stały znak.
\end{przyklad}

\section{Orientacja, współrzędne Plückera i odwzorowanie Gaussa}
\label{sec:lin29-grass-plucker}

\begin{definicja}[Grassmannian zorientowany]
\index{grassmannian zorientowany@Grassmannian zorientowany}
Punkt $\operatorname{Gr}^+_k(\R^n)$ to para $(L,o)$:
$k$-płaszczyzna $L$ wraz z wyborem jej orientacji $o$.
Zapomnienie orientacji daje dwukrotne nakrycie
$\operatorname{Gr}^+_k(\R^n)\to\operatorname{Gr}_k(\R^n)$
dla $0<k<n$: nad każdą płaszczyzną są dokładnie
dwie możliwe orientacje. W mapie wykresowej
orientację przenosimy z $P$ na $\Gamma(A)$ przez
$u\mapsto u+Au$, więc lokalnie nakrycie jest
sumą dwóch kopii tej samej mapy.
\end{definicja}

\begin{przyklad}[Płaszczyzny w $\R^3$]
Płaszczyźnie $L\in\operatorname{Gr}_2(\R^3)$
odpowiada jej nieukierunkowana prosta normalna
$L^\perp$, a zatem
$\operatorname{Gr}_2(\R^3)\cong\mathbb{RP}^2$.
Gdy $L$ ma orientację, wybierzmy dodatnią bazę
ortonormalną $(u,v)$ i połóżmy $n=u\times v$.
Inna dodatnia baza różni się obrotem w $L$,
więc daje ten sam $n$. W drugą stronę $n\in S^2$
wyznacza $L=n^\perp$ i orientację, dla której
$(u,v,n)$ jest dodatnią bazą $\R^3$. Otrzymujemy
\[
 \operatorname{Gr}^+_2(\R^3)\cong S^2,\qquad
 \operatorname{Gr}_2(\R^3)\cong S^2/(n\sim -n).
\]
Zmiana orientacji płaszczyzny zamienia $n$ z $-n$.
Dla ustalenia gładkości zauważmy, że normalna z lokalnej dodatniej
ramki ortonormalnej jest gładka. W drugą stronę rzut
na $n^\perp$ ma macierz $I-nn^T$, gładką w $n$;
lokalny wybór orientacji jest stały na każdym z dwóch płatów.
Są to więc dyfeomorfizmy, nie tylko bijekcje zbiorów.
\end{przyklad}

\begin{figure}[!htbp]
\centering
\begin{tikzpicture}[>=Stealth,line cap=round,scale=1.02]
 \filldraw[fill=manifoldblue!10,draw=manifoldblue!65!black,thick]
 (-4.6,-.52)--(-1.3,-.52)--(-.45,.08)--(-3.75,.08)--cycle;
 \draw[->,very thick,manifoldred] (-2.55,-.22)--(-2.55,1.05)
   node[above,font=\small] {$n$};
 \draw[->,thick,manifoldgold!85!black,dashed]
   (-2.55,-.22)--(-2.55,-1.49)
   node[below,font=\small] {$-n$};
 \node[font=\small,manifoldblue!80!black] at (-1.2,.35) {$L$};
 \draw[->,manifoldblue!70!black] (-3.35,-.31)--(-2.45,-.31)
   node[midway,below,font=\scriptsize] {$u$};
 \draw[->,manifoldblue!70!black] (-3.35,-.31)--(-2.85,.01)
   node[near end,above,font=\scriptsize] {$v$};
 \shade[inner color=white,outer color=manifoldblue!43] (2.1,0) circle (1.12);
 \draw[manifoldblue!75!black,thick] (2.1,0) circle (1.12);
 \draw[manifoldgreen!65,thin] (.98,0) arc[start angle=180,end angle=360,x radius=1.12,y radius=.27];
 \draw[manifoldgreen!65,thin,dashed] (3.22,0)
   arc[start angle=0,end angle=180,x radius=1.12,y radius=.27];
 \fill[manifoldred] (2.1,1.12) circle (2.3pt)
   node[above,font=\small] {$n$};
 \fill[manifoldgold!85!black] (2.1,-1.12) circle (2.3pt)
   node[below,font=\small] {$-n$};
 \draw[->,manifoldblue!75!black,thick] (-.05,.15)--(.68,.15)
   node[midway,above,font=\small] {$\perp$};
 \node[font=\small] at (2.1,-1.65) {$S^2\longrightarrow\mathbb{RP}^2$};
\end{tikzpicture}
\caption{Orientacja płaszczyzny w $\R^3$ wybiera jedną z
dwóch jednostkowych normalnych. Zapomnienie orientacji
utożsamia punkty antypodyczne $n$ i $-n$ na sferze.}
\label{fig:lin29-grass-normalne}
\end{figure}

\begin{stwierdzenie}[Bazy ortonormalne i przestrzeń ilorazowa]
\label{prop:lin29-grass-stiefel}
Niech $V_k(\R^n)$ będzie zbiorem uporządkowanych
ortonormalnych $k$-ramek $(v_1,\ldots,v_k)$.
Odwzorowanie $(v_i)\mapsto\operatorname{span}(v_i)$ ma
włókna izomorficzne z $\mathrm O(k)$.
Ponadto, jako gładkie przestrzenie,
\[
 \operatorname{Gr}_k(\R^n)
   \cong\mathrm O(n)/(\mathrm O(k)\times\mathrm O(n-k)),
\quad
 \operatorname{Gr}_k(\C^n)
   \cong\mathrm U(n)/(\mathrm U(k)\times\mathrm U(n-k)).
\]
\end{stwierdzenie}
\begin{proof}
Dwie ramki wyznaczają tę samą $k$-płaszczyznę wtedy
i tylko wtedy, gdy zmiana bazy między nimi należy do
$\mathrm O(k)$. Wykonując ortonormalizację Grama--Schmidta
na wektorach $u_i+Au_i$ z mapy wykresowej, dostajemy
gładki lokalny wybór ramki; włókna są więc lokalnie
kopiami $\mathrm O(k)$. Grupa $\mathrm O(n)$ działa
tranzytywnie na płaszczyznach: wybieramy bazy ortonormalne
dwóch płaszczyzn, uzupełniamy je do baz $\R^n$ i bierzemy
odwzorowanie przenoszące pierwszą pełną bazę na drugą.
Element zachowujący standardowe $\R^k$ zachowuje też
jego dopełnienie prostopadłe, czyli ma postać blokową
$\mathrm{diag}(B,C)$ z $B\in\mathrm O(k)$,
$C\in\mathrm O(n-k)$. To daje pierwszy iloraz.
Wyjaśnijmy gładką strukturę ilorazu. W rozkładzie ortogonalnym
$P\oplus Q$ uzupełniamy ramkę wykresu przez ramkę
$\Gamma(-A^*)\subset Q\oplus P$ jego dopełnienia prostopadłego:
jej wektory mają postać $-A^*w+w$, $w\in Q$.
Ortonormalizacja obu ramek daje macierz $s(A)\in\mathrm O(n)$
zależną gładko od $A$. Każda macierz ortogonalna nad tą mapą
ma jednoznaczną postać $s(A)\operatorname{diag}(B,C)$.
Parametr $A$ odzyskujemy z obrazu pierwszych $k$ kolumn,
a bloki $B,C$ z iloczynu $s(A)^{-1}O$. Daje to lokalny
produkt i jego gładką odwrotność; po podzieleniu przez bloki
pozostaje dokładnie mapa wykresowa grassmannianu.
Argument dla $\C$ jest taki sam, z bazami unitarnymi,
sprzężeniem hermitowskim i grupą $\mathrm U(n)$.
\end{proof}

Iloczyn zewnętrzny z Sekcji~\ref{sec:lin11-zewnetrzna}
pozwala zapisać $k$-płaszczyznę jednym wielowektorem,
określonym z dokładnością do niezerowego skalara.
Po wyborze bazy $(e_i)$ przestrzeni $V$ oraz zapisaniu
$v_j=\sum_i v_{ij}e_i$ współczynnik
$p_I$ przy $e_{i_1}\wedge\cdots\wedge e_{i_k}$
ma postać minora
$\det(v_{i_a j})_{1\leq a,j\leq k}$, gdzie
$I=(i_1<\cdots<i_k)$; jest to
wzór z Wniosku~\ref{cor:lin11-wedge-minory}.
Liczby $p_I$ nazywamy \emph{współrzędnymi Plückera}.
Zmiana bazy $L$ mnoży je wszystkie przez tę samą
niezerową liczbę $\det C$, więc geometryczną informację
niesie ich klasa rzutowa $[p_I]$.

\begin{twierdzenie}[Zanurzenie Plückera]
\label{thm:lin29-grass-plucker}
Jeśli $(v_1,\ldots,v_k)$ jest bazą $L$, to wzór
\begin{equation}\label{eq:lin29-plucker}
 \mathrm{Pl}(L):=[v_1\wedge\cdots\wedge v_k]
 \in\mathbb P(\Lambda^k V)
\end{equation}
określa gładkie zanurzenie
$\operatorname{Gr}_k(V)\hookrightarrow\mathbb P(\Lambda^k V)$.
Nawias $[\xi]$ oznacza prostą $\mathbb K\xi$,
$\xi\ne0$. Obraz składa się dokładnie z klas
niezerowych \emph{wielowektorów prostych}.
\end{twierdzenie}
\begin{proof}
Zmiana bazy $(v_i)$ macierzą $C$ mnoży iloczyn
zewnętrzny przez $\det C\ne0$ na mocy
Twierdzenia~\ref{thm:lin11-det-wedge}, więc klasa
rzutowa jest dobrze określona. Jeżeli
$\xi=v_1\wedge\cdots\wedge v_k\ne0$, to
\[
 L=\{w\in V:w\wedge\xi=0\}.
\]
Wektory z $L$ dają zero przez powtórzenie czynnika.
Jeśli $w\notin L$, rodzina $(v_1,\ldots,v_k,w)$
jest liniowo niezależna; po uzupełnieniu do bazy
Twierdzenie~\ref{thm:lin11-zew-baza} mówi, że jej
iloczyn zewnętrzny jest niezerowy. Zatem
$\mathrm{Pl}(L)$ wyznacza $L$ jednoznacznie.

Pokażmy, skąd bierze się gładki odwrotny wzór na obrazie.
W mapie $V=P\oplus Q$ wybierzmy bazy
$e_1,\ldots,e_k$ przestrzeni $P$ oraz
$e_{k+1},\ldots,e_n$ przestrzeni $Q$ i zapiszmy
$Ae_j=\sum_{a=1}^{n-k}a_{aj}e_{k+a}$.
W wielowektorze
$\xi_A=(e_1+Ae_1)\wedge\cdots\wedge(e_k+Ae_k)$
współczynnik $p_{1\cdots k}$ wynosi $1$.
Współczynnik przy
$e_1\wedge\cdots\wedge\widehat{e_j}\wedge\cdots
\wedge e_k\wedge e_{k+a}$ wynosi
$(-1)^{k-j}a_{aj}$: przesuwamy $e_{k+a}$
przez $k-j$ późniejszych czynników.
Z ilorazów tych współczynników przez $p_{1\cdots k}$
odzyskujemy wszystkie $a_{aj}$.
Jeśli dla dowolnego indeksu $I$ mamy $p_I\ne0$,
to rzut bazy $L$ na płaszczyznę
$P=\operatorname{span}\{e_i:i\in I\}$ ma wyznacznik $p_I$.
Jest więc izomorfizmem $L\to P$ i możemy stosować
analogiczny wzór w mapie $U_P$, dzieląc przez $p_I$.
Odwzorowanie~\eqref{eq:lin29-plucker} jest więc gładkie.
Wzory odzyskujące $a_{aj}$ są określone na całym otwartym
zbiorze przestrzeni rzutowej $p_{1\cdots k}\ne0$, nie tylko
na obrazie. Ich złożenie z $\mathrm{Pl}$ jest odwzorowaniem identycznościowym
w parametrach $A$; różniczkowanie daje lewą odwrotność różniczki
$\mathrm{Pl}$, która jest zatem injektywna.
Otrzymaliśmy immersję. Grupy $\mathrm O(n)$ i $\mathrm U(n)$
są zwarte jako domknięte i ograniczone zbiory macierzy.
Ich ilorazy z~\ref{prop:lin29-grass-stiefel}, a więc oba
grassmanniany, są zwarte jako ciągłe obrazy grup zwartych.
Ponieważ przestrzeń rzutowa jest Hausdorffa,
ciągła iniekcja ze zwartego grassmannianu jest
homeomorfizmem na obraz. To daje zanurzenie i domkniętość obrazu.
\end{proof}

\begin{przyklad}[Relacja Plückera dla $\operatorname{Gr}_2(\R^4)$]
Zapiszmy $\xi=\sum_{1\leq i<j\leq4}p_{ij}e_i\wedge e_j$.
Rozwińmy $\xi\wedge\xi$. Składnik z powtórzonym
indeksem znika, więc zostają tylko pary $12|34$, $13|24$
i $14|23$ oraz te same pary w odwrotnej kolejności.
Pierwsza para jest już uporządkowana; w drugiej
przestawienie $1,3,2,4$ do $1,2,3,4$ wymaga jednej
transpozycji, a w trzeciej przestawienie $1,4,2,3$
wymaga dwóch. Znaki wynoszą zatem kolejno $+,-,+$.
Każda para występuje dwukrotnie z tym samym znakiem,
bo zamiana dwóch bloków stopnia $2$ daje $(-1)^4=1$.
Otrzymujemy
\begin{equation}\label{eq:lin29-plucker-24}
 \xi\wedge\xi=
 2(p_{12}p_{34}-p_{13}p_{24}+p_{14}p_{23})
   e_1\wedge e_2\wedge e_3\wedge e_4.
\end{equation}
Jeżeli $\xi=u\wedge v$, to $\xi\wedge\xi=0$,
więc współczynniki spełniają relację w nawiasie.
Odwrotnie, dla $\xi\ne0$ z taką relacją możemy
przestawić współrzędne tak, by $p_{12}\ne0$.
Po podzieleniu przez $p_{12}$ ustawiamy w wykresie
$e_1+ae_3+ce_4$, $e_2+be_3+de_4$ liczby
$b=p_{13}/p_{12}$, $d=p_{14}/p_{12}$,
$a=-p_{23}/p_{12}$, $c=-p_{24}/p_{12}$.
Pięć współczynników się zgadza z definicji,
a szósty jest równy $ad-bc=p_{34}/p_{12}$
właśnie dzięki relacji: po podstawieniu
$ad-bc=(p_{13}p_{24}-p_{14}p_{23})/p_{12}^2$,
a relacja daje w liczniku $p_{12}p_{34}$.
To dowodzi również
wystarczalności.

Policzmy wszystkie współrzędne dla
$L=\operatorname{span}(e_1+2e_3,e_2+e_4)$:
\[
 (e_1+2e_3)\wedge(e_2+e_4)
  =e_1\wedge e_2+e_1\wedge e_4
   -2e_2\wedge e_3+2e_3\wedge e_4.
\]
Mamy więc $(p_{12},p_{13},p_{14},p_{23},p_{24},p_{34})
=(1,0,1,-2,0,2)$ i relację $1\cdot2+1\cdot(-2)=0$.
Z kolei współrzędne wykresu odzyskujemy jako
$(a,b,c,d)=(2,0,0,1)$, zgodnie z wyjściową bazą.
W tym przypadku $\dim\operatorname{Gr}_2(\R^4)=4$,
podczas gdy $\mathbb P(\Lambda^2\R^4)$ ma wymiar $5$:
jedna relacja $p_{12}p_{34}-p_{13}p_{24}+p_{14}p_{23}=0$
opisuje więc cały obraz Plückera jako kwadrykę rzutową.
Dla kontrastu
$e_1\wedge e_2+e_3\wedge e_4$ ma
$p_{12}p_{34}=1$ i pozostałe wyrazy zerowe,
więc nie jest wielowektorem prostym, zgodnie
z wcześniejszym Przykładem w
Sekcji~\ref{sec:lin11-baza-wedge}.
\end{przyklad}

\begin{twierdzenie}[Odwzorowanie Gaussa i wiązka styczna]
\label{thm:lin29-grass-gauss}
Niech $M^m\subset\R^N$ będzie gładką podrozmaitością.
Przypisanie
$G(p)=T_pM\in\operatorname{Gr}_m(\R^N)$
jest gładkim \emph{odwzorowaniem Gaussa}. Wiązki
$TM$ i normalna $\nu M$ są odpowiednio
cofnięciami $G^*\gamma^m$ i $G^*\mathcal Q$
z~\eqref{eq:lin29-grass-ciag}.
\end{twierdzenie}
\begin{proof}
W mapie podrozmaitości zapisujemy jej parametryzację
$F:U\subset\R^m\to\R^N$. Wektory
$\partial_iF(x)$ tworzą gładką bazę $T_{F(x)}M$.
W pobliżu ustalonego punktu pewien minor $m\times m$
macierzy $DF$ pozostaje niezerowy. W odpowiadającym
mu rozkładzie $\R^N=P\oplus Q$ przestrzeń
$T_{F(x)}M$ jest wykresem macierzy
$A(x)=D_QF(x)(D_PF(x))^{-1}$,
gładkiej w $x$. Stąd $G$ jest gładka.
Włókno $(G^*\gamma^m)_p$ to z definicji
$\gamma^m_{G(p)}=T_pM$; lokalne ramki po obu
stronach są tymi samymi $\partial_iF$.
Włókno $(G^*\mathcal Q)_p$ to
$\R^N/T_pM$, czyli $\nu_pM$ (po utożsamieniu przez iloczyn skalarny) z przestrzenią normalną $(T_pM)^\perp$.
Zgodność lokalnych trywializacji daje oba
izomorfizmy wiązek.
\end{proof}

\begin{przyklad}[Odwzorowanie Gaussa paraboli]
Dla paraboli $M=\{(t,t^2):t\in\R\}\subset\R^2$
weźmy parametryzację $F(t)=(t,t^2)$. Wówczas
\[
 F'(t)=(1,2t),\qquad
 G(F(t))=\R(1,2t),\qquad a(t)=2t
\]
w mapie $\operatorname{Gr}_1(\R^2)$ z
Przykładu~\ref{prz:lin29-grass-rp1-mapy}.
W punkcie $p=F(0)$ różniczka odwzorowania Gaussa
odpowiada, według~\eqref{eq:lin29-grass-styczna},
odwzorowaniu $T_pM\to\R^2/T_pM$ danemu przez
$e_1\mapsto2[e_2]$. Liczba $2$ mierzy tutaj
pierwszy rząd zmiany kierunku stycznej, podczas gdy
położenie punktu $F(t)$ nie występuje we współrzędnej $a(t)$.
Po utożsamieniu ilorazu z prostopadłą prostą
jednostkowa normalna ma postać
$(-2t,1)/\sqrt{1+4t^2}$.
\end{przyklad}

\begin{uwaga}[Dlaczego wiązka tautologiczna jest uniwersalna?]
Na zwartej rozmaitości każdą gładką rzeczywistą wiązkę $E$
rzędu $k$ można umieścić w trywialnej wiązce
$M\times\R^N$ dla pewnego $N$.
Oto konstrukcja. Wybieramy skończone pokrycie
trywializujące $U_i$, jego mniejsze podpokrycie
oraz funkcje gładkie $\beta_i$ o nośnikach w $U_i$,
takie że dla każdego $p$ co najmniej jedna
$\beta_i(p)$ jest niezerowa. Jeśli $[v]_i\in\R^k$
oznacza współrzędne $v\in E_p$ na $U_i$,
definiujemy różnowartościowe odwzorowanie liniowe w $v$
\[
 j_p(v)=\bigl(\beta_1(p)[v]_1,\ldots,
                  \beta_r(p)[v]_r\bigr)\in\R^{kr},
\]
przyjmując wartość $0$ dla bloku poza $U_i$.
Podparcie $\beta_i$ wewnątrz $U_i$ zapewnia
gładkość po przedłużeniu zerem; co najmniej
jeden niezerowy blok zapewnia iniektywność $j_p$.
Obrazy $j_p(E_p)$ zmieniają się gładko, ponieważ
lokalnie tworzą kolumny macierzy gładkiej i
stałego rzędu $k$, a więc definiują odwzorowanie
$f:M\to\operatorname{Gr}_k(\R^{kr})$.
Wprost z definicji włókien mamy
$E\cong f^*\gamma^k$. Ta konkretna konstrukcja
wyjaśnia użyteczność grassmannianów dla topologii
wiązek; nie potrzebuje jeszcze teorii klas
charakterystycznych.
W przypadku wiązki zespolonej ten sam rachunek
prowadzimy w $\C^{kr}$ i otrzymujemy mapę do
$\operatorname{Gr}_k(\C^{kr})$.
\end{uwaga}

